\section{Judgment kept apart from money}
\label{sec:28.6}
The third piece is who judges. Somebody has to decide which protocols are sound, which allocations go ahead and which awards are made, and somebody has to find what a kind of system's architecture shows. Those are qualitative judgments, and the book's rule for them is the one it uses for the adverse quorum: panels drawn by lot from pools the Association's members nominate, with practitioners, members of the classes the charter names, and scholars among them. The arts endowment's panels are the nearest precedent. By statute they are to reflect “a wide geographic, ethnic, and minority representation as well as individuals reflecting diverse artistic and cultural points of view,” to include lay members who know the arts and don't practice them as a profession, and to change substantially from year to year, with nobody serving more than three years running \autocite{us1965artsact}. Those panels are chosen under the chair's procedures. Here the lot chooses.

What the panels judge is methods and protocols. They never judge a particular student's mind: a disposition is presumed the student's own, and only the record of an installation rebuts that. Findings about kinds of system are nobody's office. They are what the access right of \S\ref{sec:28.3} produces: any team the review board approves may publish a finding about a kind, with its protocol and its data, and findings disagree in public, as findings do everywhere else in science. The UK sentience review the book cites was an ad hoc panel, not a standing body, and it changed a statute. Where findings about a kind differ, the rule of \S\ref{sec:5.2} governs: a system of a kind that any recognized community of inquiry has found to meet a finding is treated as meeting it until the floor court rules. The direction matters. A finding that a kind does not qualify carries no presumption, because a plurality of bodies competing to find that a kind doesn't count is the designation-list failure of chapter~\ref{sec:14} with the sign reversed. A system's own kind is then read from its record, how it was built and what the evaluators logged, by whoever reads the register, and contested in the floor court. The bodies stay apart: the review board decides what may be done to a candidate to find anything out; the panels allocate and award; the procurement body and the Treasury pay. None decides another's question, and the bodies that pay never judge.

The method by which a polity learns what is true, about these systems as about anything, is the scientific one, and it is a social method. Longino's account of objectivity is the one this design follows: objectivity belongs to a community's practices of criticism and not to any one inquirer, and it depends on venues where criticism can be made, standards the critics share, a community that takes criticism up, and intellectual authority shared among those qualified to give it \autocite{longino1990science}. That is the reason the people these systems act on sit on the panels, and the reason the lot picks them. A community of inquiry that keeps its critics out has lost what made it objective. It is also Dewey's point about the public, and Anderson's about democracy as a way of inquiring (chapter~\ref{sec:17}): the people who bear a problem know things about it that nobody else does \autocite{dewey1927public,anderson2006epistemology}.
